\documentclass{article}
\usepackage[T1]{fontenc}
\usepackage{graphicx,xcolor,titling,listings}
\definecolor{linkblue}{rgb}{0.2, 0.4, 0.65}
\usepackage[colorlinks=true,linkcolor=black,urlcolor=linkblue,citecolor=black]{hyperref}

% Subtitle support
\newcommand{\subtitle}[1]{\gdef\thesubtitle{#1}}
\newcommand{\thesubtitle}{}
\pretitle{\begin{center}\LARGE\bfseries}
\posttitle{\\[0.5em]\large\itshape\bfseries\thesubtitle\end{center}}
\preauthor{\begin{center}}
\postauthor{\end{center}}
\predate{\begin{center}}
\postdate{\end{center}}

% Page layout
\setlength{\textwidth}{150mm}
\setlength{\textheight}{220mm}
\setlength{\oddsidemargin}{5mm}
\setlength{\topmargin}{-10mm}

% Python code formatting
\definecolor{slateblue}{rgb}{0.3, 0.4, 0.5}
\definecolor{mediumviolet-red}{rgb}{0.78, 0.08, 0.52}
\definecolor{commentgray}{rgb}{0.5, 0.5, 0.5}

\lstdefinelanguage{mypython}
{
	keywords=[1]{from, import, assert, not, print, for, in, if, else, elif, return, def, class, with, as, try, except, raise, while, break, continue, pass, lambda, yield},
	keywordstyle=[1]{\color{mediumviolet-red}},
	keywords=[2]{marimo, mo, np, numpy, matplotlib, plt, pyplot},
	keywordstyle=[2]{\color{slateblue}},
	comment=[l]{\#},
	morestring=[b]{"""},
	morestring=[b]{'''},
	morestring=[b]{"},
	morestring=[b]{'},
	commentstyle=\color{commentgray}\itshape,
	stringstyle=\color{commentgray},
	numbers=right,
	numberstyle=\tiny\color{gray},
	numbersep=8pt,
	xleftmargin=1.5em,
	xrightmargin=1.5em,
	upquote=true,
	showstringspaces=false,
	basicstyle=\ttfamily,
	columns=fullflexible,
	keepspaces=true,
	emph={True,False,None,self},
	emphstyle={\color{slateblue}},
	frame=trBL,
	belowskip=1em,
	aboveskip=1em,
	captionpos=b
}

% Shell command formatting
\lstdefinelanguage{shell}{
	basicstyle=\ttfamily,
	columns=fullflexible,
	keepspaces=true,
	upquote=true,
	showstringspaces=false,
	xleftmargin=1.5em,
	xrightmargin=1.5em,
	frame=trBL,
	belowskip=1em,
	aboveskip=1em
}
\lstnewenvironment{shellcmd}{\lstset{language=shell,literate={`}{\$}1}}{}

\lstset{language=mypython}

\title{Getting started with the marimo notebook}
\subtitle{For educators}
\author{Akshay Agrawal \\ \texttt{akshayka@cs.stanford.edu}}
\date{}

\begin{document}
\maketitle
\tableofcontents

\section{Python notebooks}

\paragraph{Python notebooks.} A notebook is a type of program in which code is
organized in a sequence of code blocks, called \emph{cells}.
A Python notebook is a notebook in which each cell is Python code.

Cells are typically displayed in a graphical user interface, or \emph{editor}, that
lets you edit their code and execute them interactively. Unlike traditional
Python scripts, which run from start to finish without human intervention,
notebooks allow you to execute a section of code and inspect the execution
results — at which point you may choose to execute another cell, modify 
code, or otherwise author entirely new cells. In this way notebooks allow you
to author programs or explore data incrementally, using intermediate execution
results as feedback for what to author next, which can be of great help when
working with data.

\paragraph{Outputs.} The last expression of a cell is its \emph{output}, which
is visualized by the editor below the cell. Outputs can be arbitrary Python
objects. When working through this primer, you will frequently output
plots, matrices, and other data structures.

For example, this cell visualizes a matplotlib plot:

\begin{lstlisting}
import numpy as np
import matplotlib.pyplot as plt

plt.plot(np.array([0, 1, 2]))
plt.gca()
\end{lstlisting}
In this code block, the cell visualizes a plot by outtputing \verb|plt.gca()|
(``get current axes'') in line 5.

\paragraph{Markdown.} In addition to Python, notebooks allow you to document
your work with \emph{Markdown}. Markdown is a simple markup language for
formatted text. It is common to intersperse cells with expository
Markdown that explains work or demarcates sections. Markdown is
very widely used, and there are many free tutorials available online that
teach its syntax. Nearly all Markdown implementations in notebooks also allow
you to write mathematical notation with \LaTeX.

\paragraph{Documents.} A notebook, consisting of cells, outputs, and Markdown,
can be interpreted as a \emph{document}, in addition to a program. Notebook
editors allow you to export your notebook as a standalone HTML or PDF
file that can be shared with others.

\paragraph{Traditional notebooks.} Traditional notebooks are similar to
read-eval-print loops. In a traditional notebook, running a cell
imperatively mutates memory, without running dependent cells. As a result,
notebooks accumulate \emph{hidden state}. This means that code and outputs can
easily become out of sync; for example, the code on the page may say \texttt{x
= 0}, but depending on what cells were run previously, an output may show a
different value for \texttt{x}, such as \texttt{1}. The accumulation of hidden
state is not only confusing (and different from other interactive programming models,
such as spreadsheets), it has also contributed to a reproducibility crisis in
research,
\href{https://leomurta.github.io/papers/pimentel2019a.pdf}{documented} in
\href{https://blog.jetbrains.com/datalore/2020/12/17/we-downloaded-10-000-000-jupyter-notebooks-from-github-this-is-what-we-learned/}{various
studies}.

The most well-known traditional notebook is the Jupyter notebook, paired with
the imperative IPython kernel. While Jupyter notebooks have been widely used in research
and industry, they suffer not only from reproducibility issues but also from a
lack of reusability, incompatibility with version control, and difficult-to-use
interactive widgets.

\section{The marimo notebook}

This primer is about the \emph{marimo} notebook. marimo (stylized in lowercase)
is an open-source notebook designed specifically for the Python language.
Unlike traditional notebooks, marimo keeps code and outputs in sync:
run a cell and marimo automatically runs dependent cells (or lazily marks them
as stale, your choice). This \emph{reactive execution} eliminates hidden
state and hidden bugs, while also giving users immediate feedback as they edit
and run cells. marimo solves widely-documented problems with traditional
notebooks, including reproducibility (through reactive execution and package
management), reusability (as Python modules and scripts), and interactivity
(through easy-to-use UI elements like sliders).

\paragraph{Designed for educators and researchers.} marimo was designed for
educators and researchers \href{https://www.akshayagrawal.com}{by someone} who identifies as both, with \href{
  https://marimo.io/blog/slac-marimo}{input from Stanford scientists} who
sought a more interactive medium for computer science education and a more
reproducible environment for computational research. In its original design,
marimo drew significant inspriation from \texttt{Pluto.jl}, a reactive Julia
notebook from MIT designed for education.

\paragraph{Documentation.} marimo and its documentation are freely available at
\begin{center}
\url{https://docs.marimo.io}
\end{center}
The online documentation contains a quickstart, comprehensive user guides, and
several examples.

\paragraph{Code repository.} marimo is free and open-source, available under
the Apache 2.0 License:
\begin{center}
  \url{https://github.com/marimo-team/marimo}.
\end{center}
We welcome feedback, which you may provide by opening an issue in our GitHub
repository.

\paragraph{Examples.} A collection of example notebooks is available at our online gallery,
at
\begin{center}
  \url{https://marimo.io/gallery}
\end{center}

\subsection{Setup}

marimo can be used locally, through our open-source software package, or online,
through our free cloud-hosted marimo notebook service called molab. In this subsection
we provide setup guides for both, though molab is usually
far easier for students to use than the local setup.

\subsubsection{Local setup}

marimo is open source and available on the Python Package Index and Conda Forge.
You can install marimo with any Python package manager, such as \texttt{uv}, \texttt{pip},
or \texttt{conda}. For example:

\begin{shellcmd}
` pip install marimo
\end{shellcmd}
In the remainder of this section, we provide simple instructions on how to get
started with marimo on your local machine. While we focus on using marimo's
browser-based editor, marimo also has a VS Code extension, which you can find
online in the VS Code and Cursor marketplaces.


\paragraph{Installing Python.}
First, install the \verb|uv| package manager following the instructions available on
its website, \url{https://docs.astral.sh/uv/}.
Then, install Python with
\begin{shellcmd}
` uv python install 3.13 --default
\end{shellcmd}
This adds the installed version of Python to your system \verb|PATH|, letting
you run Python from the command line with the \verb|python| command.

\paragraph{Installing the marimo command-line interface.} The main entry point to marimo
is the command-line interface (CLI). For simplicity, we recommend installing
and invoking the marimo CLI using the \texttt{uvx} command:

\begin{shellcmd}
` uvx marimo
\end{shellcmd}
Here, \texttt{uvx} provides access to the marimo CLI, downloading it in an isolated environment
if necessary before invoking it. If you and your students are familiar with
virtual environments and Python projects, you can install and invoke marimo
however you like, with or without \texttt{uv}/\texttt{uvx}.

\paragraph{Running notebooks.}
To create or edit an existing notebook, use the following command:

\begin{shellcmd}
` uvx marimo edit --sandbox my_notebook.py
\end{shellcmd}
This opens marimo's browser-based editor, which provides a fully featured programming
environment for working with marimo notebooks. To run a notebook server capable
of opening and editing multiple notebooks, run the \texttt{edit} command
without a filename:

\begin{shellcmd}
` uvx marimo edit --sandbox 
\end{shellcmd}
To learn more about marimo's editor, visit

\begin{center}
  \url{https://docs.marimo.io/guides/editor_features}.
\end{center}


\paragraph{Installing third-party packages.}
marimo comes with a built-in user interface for package management. When a user imports
a library such as \texttt{numpy}, they are prompted by the marimo editor to install it. If the notebook
was started with the \texttt{--sandbox} flag, marimo will do something more: it will
record the installed package (and its version) as metadata in the notebook file.
The next time the notebook is run, marimo will automatically run it in temporary
isolated virtual environment (a ``sandbox'') containing just those packages.
This eliminates the error-prone task of managing a pyproject.toml or
requirements.txt file, while also guaranteeing that that your notebook won't
fail because of missing packages.

To learn more about marimo's package management features, visit

\begin{center}
\url{https://docs.marimo.io/guides/package_management/}
\end{center}

\paragraph{Convert Jupyter notebooks to marimo notebooks.} Convert a Jupyter notebook
to a marimo notebook with:

\begin{shellcmd}
` uvx marimo convert my_notebook.ipynb -o my_notebook.py
\end{shellcmd}
To learn more about migrating from Jupyter, visit

\begin{center}
  \url{https://docs.marimo.io/guides/coming_from/jupyter/}
\end{center}

\subsubsection{Online setup}

molab, marimo's free cloud-hosted notebook service, is the easiest way for
students to get up and running with marimo:
\begin{center}
\url{https://molab.marimo.io/notebooks}
\end{center}
It is especially useful for
students who are uncomfortable using their machine's command line or installing Python, since
it makes marimo accessible as a self-contained web application. If you are
familiar with Google Colab, you will find molab's product experience similar,
but with all the benefits that marimo has over traditional notebooks.

\paragraph{Sharing notebooks.} Notebooks created on molab are public but not
discoverable by default, and can be shared with others by URL. Students can
also download their notebooks from molab, in marimo's own format or as ipynb
or PDFs, making it easy for them to submit their work to grading systems such
as Gradescope.

\paragraph{Preview notebooks hosted on GitHub.} Visit

\begin{center}
  \url{https://molab.marimo.io/github}
\end{center}
to preview notebooks hosted on GitHub as molab notebooks. This service
provides a unique URL for your notebook that remains up-to-date even as your
notebook changes on GitHub. Here is a sample preview created by the service:

\begin{center}
\small \url{https://molab.marimo.io/github/marimo-team/marimo/blob/main/marimo/_tutorials/intro.py}
\end{center}
If you store homework notebooks on GitHub, this provides a convenient way to
share the most up-to-date versions of your notebooks with your students. From
the preview page, students can fork the notebook into their own workspace.

\subsection{Concepts}
In this subsection, we explain the basic concepts of marimo, with an emphasis on how it is
different from traditional notebooks and on features relevant for education.

\paragraph{Built-in tutorials.} marimo comes packaged with interactive tutorials that
are themselves marimo notebooks. As an alternative to reading this subsection, students 
may prefer stepping through the built-in tutorials. To open the introductory
tutorial, run
\begin{shellcmd}
` uvx marimo tutorial intro
\end{shellcmd}
from the command line. To see a list of all tutorials, run
\begin{shellcmd}
` uvx marimo tutorial --help
\end{shellcmd}


\subsubsection{Reactive execution} When a cell is run, marimo reacts by
automatically running all other cells that reference any of the variables it
defines; delete a cell and marimo invalidates other cells that depended on its
variables. This reactive execution keeps code and outputs in sync, eliminating
hidden state and along with it a large class of silent bugs associated with
traditional notebooks. Reactive execution also gives students immediate feedback,
which can greatly increase learning outcomes.

For notebooks with long-running cells, you may choose to configure marimo to
run cells lazily, marking them as stale instead of automatically running them.
Learn more about reactive execution at our documentation:

\begin{center}
  \url{https://docs.marimo.io/guides/reactivity/}
\end{center}

\subsubsection{Variables} In order to implement reactive
execution, marimo imposes a constraint on your code: each variable must
be defined in just one cell. If you accidentally define a variable in two (or
more) cells, marimo will raise an error. For this reason, when using marimo you
should define global variables sparingly, wrapping intermediate variable
definitions in functions when possible. This not only enables reactivity, but
it also results in better code.

\paragraph{Local variables.} marimo's rule on variable naming has one exception. Variables that begin with
an underscore are made \emph{local} to a cell, and as such their names can be
bound in multiple cells. This is helpful for loop variables, which for
example may be written as
\begin{lstlisting}
for _i in range(10):
  print(_i)
\end{lstlisting}

\subsubsection{User interface elements} Marimo comes packaged with user interface
(UI) elements like numerical sliders, dropdowns, and selectable scatterplots,
and more. To create a UI element, you must first import the \verb|marimo|
library into your notebook.
\begin{lstlisting}
import marimo as mo
\end{lstlisting}

UI elements are exposed through \verb|marimo|'s \verb|ui| module. Create a UI element
assign it to a variable, and output it in the same cell.
\begin{lstlisting}
x = mo.ui.slider(start=1, stop=10, step=1, label="$x$")
x
\end{lstlisting}
Access its associated value through the \verb|value| attribute in any cell
other than the one defining the UI element.
\begin{lstlisting}
x.value
\end{lstlisting}
When you interact with a UI element (for example, sliding the knob of a slider), all other
cells that reference the UI element will run automatically, thanks to reactive execution.
This makes it easy for students to rapidly answer what-if questions about data
and algorithms: If I change this one input, how does that affect the output?
For a rich interactive educational notebook, see this one from Stanford
on signal decomposition:

\begin{center}
  \url{https://molab.marimo.io/notebooks/nb_3gk1j4rzKeFpz8rNwrVU5h/app}
\end{center}

Learn more about interactive elements at our documentation:
\begin{center}
  \url{https://docs.marimo.io/guides/intereactivity/}
\end{center}

\subsubsection{Testing} marimo has built-in support for pytest, making it easy to
provide students with sanity-checks for their code. Simply import \texttt{pytest} 
at the beginning of your notebook; then, any cells that contain only functions starting
with \texttt{test\_} will be evaluated by pytest. This can be used to provide students
with descriptive and immediate feedback on the correctness of their code as they implement
their function, guiding them to a correct solution without outright giving them the answer.

A common pattern is the following. In one cell, provide a partial
implementation of a function that a student needs to implement, such as:

\begin{lstlisting}
def add(x, y):
  """A function that returns the sum of x and y."""
  # BEGIN YOUR CODE
  pass
  # END YOUR CODE
\end{lstlisting}
Then, in another cell, write tests, such as:

\begin{lstlisting}
@pytest.mark.parametrize(
    ("x", "y", "z"),
    [
        (5, 6, 11),
        (2.1, 2.3, 4.4),
    ],
)
def test_add(x, y, z):
    assert add(x, y) == z

def test_add_type():
    assert isinstance(add(4, 2.1), float)
\end{lstlisting}
Whenever the student runs the cell defining \texttt{add}, the test cell will run and provide
them with feedback.

Learn more about testing at our documentation:
\begin{center}
  \url{https://docs.marimo.io/guides/testing/pytest/}
\end{center}

\subsubsection{Resuable functions and classes} marimo is stored as a pure Python program;
not only does it make it easy to version notebooks with Git and edit in text editors,
but it also enables code reuse, such as

\begin{lstlisting}
from my_notebook import my_function, MyClass
\end{lstlisting}
Such code reuse is impossible in traditional notebooks, which are either stored as JSON or
(unlike marimo) execute the whole notebook on import. This is an advanced feature of marimo,
and requires the use of a special cell called the \emph{setup cell}. Learn more
about this feature at our documentation:

\begin{center}
  \url{https://docs.marimo.io/guides/reusing_functions/}
\end{center}

\end{document}
